\chapter*{Preface}




\subsection*{Causal inference research and education in the past decade}



The past decade has witnessed an explosion of interest in research and education in causal inference, due to its wide applications in biomedical research, social sciences, tech companies, etc. It was quite different even ten years ago when I was a Ph.D. student in statistics. At that time, causal inference was not a mainstream research topic in statistics and very few undergraduate and graduate programs offered courses in causal inference.  In the academic world of statistics, many people were still very skeptical about the foundation of causal inference.  Many leading statisticians were reluctant to accept causal inference because of the fundamental conceptual difficulties, which differ from the traditional training of mathematical statistics. 


The applications of causal inference in empirical research have changed the field of statistics in both research and education. In the end, statistics is not only about abstract theory but also about solving real-world problems.  Many talented researchers have joined the effort to advance our knowledge of causal inference.  Many students are eager to learn state-of-the-art theory and methods in causal inference so that they are better equipped to solve problems from various fields. 



Due to the needs of the students, my colleagues encouraged me to develop a course in causal inference. Initially, I taught a graduate-level course cross-listed under Political Science and Statistics, which was taught by my former colleague Jas Sekhon for many years at UC Berkeley. Later, I developed this course for both undergraduate and graduate students. At UC Berkeley, the course numbers for ``Causal Inference'' are Stat 156 and Stat 256, with undergraduate students in Stat 156 and graduate students in Stat 256.  Students in both sessions used the same lecture notes and attended the same lectures given by me and my teaching assistants, although they needed to finish different homework problems, reading assignments, and final projects.  


Given the mixed levels of technical preparations of my students, the most challenging part of my teaching was to balance the interests of both undergraduate and graduate students. On the one hand, I wanted to present the materials in an intuitive way and only required the undergraduate students to have the basic knowledge of probability, statistics, linear regression, and logistic regression. On the other hand, I also wanted to introduce recent research topics and results to the graduate students. This book is a product of my efforts in the past seven years.  



\subsection*{Recommendations for instructors}



This book contains 29 chapters in the main text and 3 chapters in the appendix.  UC Berkeley is on the semester system and each semester has 14 weeks of lectures. I could not finish all 32 chapters in one semester.  Here are some recommendations based on my own teaching experience.


\paragraph*{Appendix}
I started with the chapters in the main text but asked my teaching assistants to review the basics in Chapters \ref{chapter::basic-prob-append} and \ref{appendix::basic-linear-regression}.  To encourage the students to review Chapters \ref{chapter::basic-prob-append}--\ref{chapter::SRS-appendix} before reading the main text, I also assigned several homework problems from Chapters \ref{chapter::basic-prob-append}--\ref{chapter::SRS-appendix} at the beginning of the semester.



\paragraph*{Part \ref{part::introduction}}
The key topic in Chapter \ref{chapter::correlationassociation} is the Yule--Simpson Paradox.  Chapter \ref{chapter::potential-outcomes} introduces the notion of potential outcomes, which is the foundation for the whole book. 



\paragraph*{Part \ref{part::rcts}}
Different researchers and instructors may have quite different views on the materials in Part \ref{part::rcts} on randomized experiments.  I have talked to many friends about the pros and cons of the current presentation in Part \ref{part::rcts}.  Causal inference in randomized experiments is relatively straightforward because randomization eliminates unmeasured confounding. So some friends feel that Chapters \ref{ch::frt-cre}--\ref{chapter::bridging} are too long at least for beginners of causal inference.  This also disappointed some of my students when I spent a month on randomized experiments. 
On the other hand,  I was trained from the book of \citet{imbens2015causal} and believed that to understand observational studies, it is better to understand randomized experiments first. Moreover, I am a big fan of the canonical research of \citet{Neyman:1923} and \citet{Fisher:1935}.  Therefore, Part \ref{part::rcts} deeply reflects my own intellectual history and personal taste in statistics. Other instructors may not want to spend a month on randomized experiments and can cover Chapters \ref{chapter::stratification-poststratification}, \ref{chapter::mpe}, \ref{chapter::unification-fisher-neyman} and \ref{chapter::bridging} quickly.


\paragraph*{Part \ref{part::observational-studies}}
Part \ref{part::observational-studies} covers the key ideas in observational studies without unmeasured confounding.  Four pillars of observational studies are outcome regression, inverse propensity score weighting, doubly robust, and matching estimators, which are covered in Chapters \ref{chapter::observational-studies}, \ref{chapter::pscore-key}, \ref{chapter::doubly-robust} and \ref{chapter::matching-obs}, respectively. 
Chapters \ref{chapter::ATT-and-other} and \ref{chapter::pscore-unification} are optional in teaching. But the results in Chapters \ref{chapter::ATT-and-other} and \ref{chapter::pscore-unification} are not uninteresting, so I sometimes covered one or two results there, asked the teaching assistants to cover more in the lab sessions, and encouraged students to read them by assigning some homework problems from those chapters. 


\paragraph*{Part \ref{part::challenges-os}}
Part \ref{part::challenges-os} is a novel treatment of the fundamental difficulties of observational studies including unmeasured confounding and overlap. However,  this part is far from perfect due to the complexities and subtleties of the issues.  Chapters \ref{chapter::evalue}, \ref{chapter::sensitivity-ACE} and \ref{chapter::overlap} are central, whereas Chapters \ref{chapter::difficulties-observational-studies} and \ref{sec::rosenbaum-sensitivity-analysis} are optional. 


\paragraph*{Part \ref{part::instrumentalvariables}}
Part \ref{part::instrumentalvariables} discusses the idea of the instrumental variable.  Chapters \ref{chapter::iv-experiment}, \ref{chapter::iv-econometric} and \ref{chapter::iv-frdd} are key, whereas Chapters \ref{chapter::iv-inequalities} and \ref{chapter::iv-mr} are optional.


\paragraph*{Part \ref{part::post-treatmentvariable}}
Part \ref{part::post-treatmentvariable} are some special topics.  They are all optional in some sense. Probably it is worth teaching Chapter \ref{chapter::mediation} given the popularity of the Baron--Kenny method in mediation analysis. 


\paragraph*{Omitted topics}
This book does not cover some popular econometric methods including the difference in differences, panel data, and synthetic controls. Instructors can use \citet{angrist2008mostly} as a reference for those topics. 
This book assumes minimal preparation for the background knowledge in probability and statistics.  Since most introductory statistics courses use the frequentists' view that assumes the unknown parameters are fixed, I adopt this view in this book and omit the Bayesian view for causal inference. In fact, many fundamental ideas of causal inference are from the Bayesian view, starting from \citet{rubin1978bayesian}. If readers and students are interested in Bayesian causal inference, please read the review paper by \citet{li2023bayesian}. 



\paragraph*{Help from teaching assistants}
My teaching assistants offered invaluable help for my courses at UC Berkeley. Since I could not cover everything in my notes, I consistently relied on them to cover some technical details or \ri{R} program sessions in their labs. 
 


\paragraph*{Solution to some homework problems}
I have also prepared the solutions to most theory problems. If you are an instructor for a causal inference course, please contact me for the solutions with detailed information about your course.  





\subsection*{Additional recommendations for readers and students}



Readers and students can first read my recommendations for instructors above.  In addition, I have two other recommendations. 


\paragraph*{Homework problems}
Each chapter of this book contains homework problems. To deepen the understanding,  it is important to try some homework problems. Moreover, some homework problems contain useful theoretical results. Even if you do not have time to figure out the details for those problems,  it is helpful to at least read the statements of the problems. 

\paragraph*{Recommended reading}
Each chapter of this book contains recommended reading.  If you want to do research in causal inference, those recommended papers can be useful background knowledge of the literature.  When I taught the graduate-level causal inference course at UC Berkeley, I assigned the following papers to the students as weekly reading from week one to the end of the semester:  
\begin{itemize}
\item 
\citet{bickel1975sex}; 
\item 
\citet{holland1986statistics};
\item
\citet{miratrix:2013};
\item
\citet{lin2013};
\item 
\citet{li2018asymptotic};
\item 
\citet{rosenbaum1983central};
\item 
\citet{lunceford2004stratification};
\item
\citet{ding2016sensitivity};
\item
\citet{pearl1995causal};
\item
\citet{angrist1996identification};
\item
\citet{imbens2014instrumental};
\item
\citet{frangakis2002principal}. 
\end{itemize}
Many students gave me positive feedback about their experience of reading the papers above. 
I recommend reading the above papers even if you do not read this book. 




 


\subsection*{Features of the book}


There are already many excellent causal inference books published in the last decade.  Some of them have profound influences on me. 
When I was in college, I read some draft chapters of \citet{imbens2015causal} from the internet. They completely challenged my way of thinking about statistics and helped to build my research interest in causal inference.  
I read \citet{angrist2008mostly} many times and have gained new insights each time I re-read it. 
\citet{rosenbaum2002design}, \citet{morgan2015counterfactuals}, and \citet{hernan2020causal} are another three excellent books from leading researchers in causal inference.
When I was preparing for the book,  \citet{cunningham2021causal}, \citet{huntington2022effect},  \citet{brumback2022fundamentals} and \citet{huber2023causal} appeared as four recent excellent books on causal inference. 


Thanks to my teaching experience at UC Berkeley, this book has the following features that instructors, students, and readers may find attractive. 
\begin{itemize}
\item
This book assumes minimal preparation for causal inference and reviews the basic probability and statistics knowledge in the appendix.

\item
This book covers causal inference from the statistics,  biostatistics, and econometrics perspectives, and draws applications from various fields. 

\item
This book uses \ri{R} code and data analysis to illustrate the ideas of causal inference. All the \ri{R} code and datasets are publicly available at Harvard Dataverse: \url{https://doi.org/10.7910/DVN/ZX3VEV}

\item
This book contains homework problems and can be used as a textbook for both undergraduate and graduate students. Instructors can also ask me for solutions to some homework problems. 
\end{itemize}




\subsection*{Acknowledgments}



Professor Zhi Geng at Peking University introduced me to the area of causal inference when I was studying in college. Professors Luke Miratrix, Tirthankar Dasgupta, and Don Rubin served on my Ph.D. thesis committee at the Harvard Statistics Department. Professor Tyler VanderWeele supervised me as a postdoctoral researcher in Epidemiology at the Harvard T. H. Chan School of Public Health.   


My colleagues at the Berkeley Statistics Department have created a critical and productive research environment.  Bin Yu and Jas Sekhon have been very supportive since I was a junior faculty.  My department chairs, Deb Nolan, Sandrine Dudoit, and Haiyan Huang, encouraged me to develop the ``Causal Inference'' course.  It has been a rewarding experience for me.  


I have been lucky to work with many collaborators,  in particular, Avi Feller,  Laura Forastiere, Zhichao Jiang, Fang Han, Fan Li, Xinran Li, Alessandra Mattei, Fabrizia Mealli, Shu Yang, and Anqi Zhao.  I will report in this book what I have learned from them. 


Many students at UC Berkeley made critical and constructive comments on early versions of my lecture notes.  As teaching assistants for my ``Causal Inference'' course,  Emily Flanagan and Sizhu Lu read early versions of my book carefully and helped me to improve the book a lot. 


Professor Joe Blitzstein read an early version of the book carefully and made very detailed comments. Addressing his comments leads to significant improvement in the book.  Professors Hongyuan Cao and Zhichao Jiang taught ``Causal Inference'' courses based on an early version of the book. They made very valuable suggestions. 


I am also very grateful for the suggestions from Young Woong Min, Fangzhou Su, Chaoran Yu, and Lo-Hua Yuan.


If you identify any errors, please feel free to email me.
